\documentclass[12pt]{article}

\usepackage[margin=50pt]{geometry}
\usepackage{amsmath}
\usepackage{amssymb}
\usepackage{graphicx}
\usepackage{xcolor}

\usepackage{titling}
\newcommand{\subtitle}[1]{%
  \posttitle{%
    \par\end{center}
    \begin{center}\large#1\end{center}
    \vskip0.5em}%
}

\graphicspath{/images}
\geometry{footskip=0.8cm}

\title{Linear Programming}
\subtitle{Graded Assignments 3}
\author{Thomas Vroom (6291496)}

\begin{document}
\maketitle

%
% Exercise 1
%
\section{Algebraic Construction of the Simplex Tableau}
Consider the following linear program:
\begin{align*}
  &\mathrm{max} \quad Z = 2x_1 - x_2 + 3x_3 \\[2pt]
  &\mathrm{s.t.} \quad 
  \begin{array}[t]{r c r c r c r c r c r c r c}
      3x_1 & + & x_2 & + & x_3 & + & \textcolor{red}{x_4} & & & & & = & 60 \\
      x_1 & - & x_2 & + & 2x_3 & & & + & \textcolor{red}{x_5} & & & = & 10 \\
      -x_1 & + & x_2 & - & x_3 & & & & & + & \textcolor{red}{x_6} & = & 20
  \end{array}\\
  &\mathrm{and} \quad x_1, x_2, x_3 \geq 0
\end{align*}

\noindent
From this LP we can extract the following information:

\begin{equation*}
  A = \begin{bmatrix}
     3 &  1 &  1 \\
     1 & -1 &  2 \\
    -1 &  1 & -1
  \end{bmatrix},\ 
  %
  b = \begin{pmatrix}
    60 \\
    10 \\
    20
  \end{pmatrix},\ 
  %
  c = \begin{pmatrix}
    2 & -1 & 3
  \end{pmatrix}
\end{equation*}

\noindent
After an arbitrary number of pivots the tableau of this LP looks like this:

\begin{table}[h]
  \centering
  \begin{tabular}{c|c|ccc|ccc|c}
  Basis & $Z$ & $x_1$ & $x_2$ & $x_3$ & $x_4$ & $x_5$ & $x_6$ & RHS \\ \hline
  $Z$ & 1 & ? & ? & -0.25 & ? & ? & ? & 27.5 \\ \hline
  ? & 0 & 0 & ? & ? & 0.25 & ? & ? & ? \\
  ? & 0 & ? & ? & ? & ? & ? & ? & ? \\
  $x_6$ & 0 & ? & ? & ? & 0 & ? & ? & 30
  \end{tabular}
\end{table}

\noindent
From this tableau we can make the following observations:

\begin{itemize}
  \item $x_6$ is basic with a value of 30.
  \item $x_3$ and $x_4$ are non-basic (they don't have an identity pattern in their column). This means $x_3 = 0$ and $x_4 = 0$.
\end{itemize}

\noindent
If we fill the constraints with these values we get the following equations:
\[ 3x_1 + x_2 = 60 \]
\[ x_1 - x_2 + x_5 = 10 \]
\[ -x_1 + x_2 + 30 = 20 \]

\newpage
\noindent
These equations are solvable using basic algebra.
This gives us the following values for all variables:
\[ x_1 = 17.5,\ x_2 = 7.5,\ x_3 = 0,\ x_4 = 0,\ x_5 = 0,\ x_6 = 30 \]

\noindent
From this we can conclude $x_1$, $x_2$ and $x_6$ are basic variables, which means they should have an identity pattern in the column of the tableau.
We know $x_6$ should have a 1 in the bottom row.
$x_1$ cannot have a 1 in the top row, which means it must have a 1 in the middle row.
This leaves $x_2$ with a 1 in the top row:

\begin{table}[h]
  \centering
  \begin{tabular}{c|c|ccc|ccc|c}
  Basis & $Z$ & $x_1$ & $x_2$ & $x_3$ & $x_4$ & $x_5$ & $x_6$ & RHS \\ \hline
  $Z$ & 1 & \textcolor{blue}{0} & \textcolor{blue}{0} & -0.25 & ? & ? & \textcolor{blue}{0} & 27.5 \\ \hline
  $\textcolor{blue}{x_2}$ & 0 & 0 & \textcolor{blue}{1} & ? & 0.25 & ? & \textcolor{blue}{0} & \textcolor{blue}{7.5} \\
  $\textcolor{blue}{x_1}$ & 0 & \textcolor{blue}{1} & \textcolor{blue}{0} & ? & ? & ? & \textcolor{blue}{0} & \textcolor{blue}{17.5} \\
  $x_6$ & 0 & \textcolor{blue}{0} & \textcolor{blue}{0} & ? & 0 & ? & \textcolor{blue}{1} & 30
  \end{tabular}
\end{table}

\noindent
From this tableau and the original LP we can make the following observations:

\begin{equation*}
    c_B = \begin{pmatrix}
      -1 & 2 & 0
    \end{pmatrix},\ 
    %
    B = \begin{bmatrix}
       1 &  3 & 0 \\
      -1 &  1 & 0 \\
       1 & -1 & 1
    \end{bmatrix}
\end{equation*}

\noindent
Inverting $B$ gives us the $B^{-1}$ matrix:

\begin{equation*}
  B^{-1} = \begin{bmatrix}
    0.25 & -0.75 & 0 \\
    0.25 &  0.25 & 0 \\
       0 &     1 & 1
  \end{bmatrix}
\end{equation*}

\noindent
From here, we can calculate $B^{-1}A$ and $c_BB^{-1}$ to complete the tableau:

\begin{equation*}
  B^{-1}A = \begin{bmatrix}
    0 & 1 & -1.25 \\
    1 & 0 & 0.75 \\
    0 & 0 & 1
  \end{bmatrix}
\end{equation*}
\begin{equation*}
  c_BB^{-1} = \begin{pmatrix}
    0.25 & 1.25 & 0
  \end{pmatrix}
\end{equation*}

\begin{table}[h]
  \centering
  \begin{tabular}{c|c|ccc|ccc|c}
  Basis & $Z$ & $x_1$ & $x_2$ & $x_3$ & $x_4$ & $x_5$ & $x_6$ & RHS \\ \hline
  $Z$ & 1 & \textcolor{blue}{0} & \textcolor{blue}{0} & -0.25 & \textcolor{red}{0.25} & \textcolor{red}{1.25} & \textcolor{blue}{0} & 27.5 \\ \hline
  $\textcolor{blue}{x_2}$ & 0 & 0 & \textcolor{blue}{1} & \textcolor{red}{-1.25} & 0.25 & \textcolor{red}{-0.75} & \textcolor{blue}{0} & \textcolor{blue}{7.5} \\
  $\textcolor{blue}{x_1}$ & 0 & \textcolor{blue}{1} & \textcolor{blue}{0} & \textcolor{red}{0.75} & \textcolor{red}{0.25} & \textcolor{red}{0.25} & \textcolor{blue}{0} & \textcolor{blue}{17.5} \\
  $x_6$ & 0 & \textcolor{blue}{0} & \textcolor{blue}{0} & \textcolor{red}{1} & 0 & \textcolor{red}{1} & \textcolor{blue}{1} & 30
  \end{tabular}
\end{table}

%
% Exercise 2
%
\newpage
\section{Mechanical Construction of the Dual}
\subsection{A. Construct the dual of the LP shown in exercise 1.}
\begin{align*}
  &\mathrm{max} \quad Z = 2x_1 - x_2 + 3x_3 \\[2pt]
  &\mathrm{s.t.} \quad 
  \begin{array}[t]{r c r c r c r}
      3x_1 & + & x_2 & + & x_3 & \leq & 60 \\
      x_1 & - & x_2 & + & 2x_3 & \leq & 10 \\
      -x_1 & + & x_2 & - & x_3 & \leq & 20
  \end{array}\\
  &\mathrm{and} \quad x_1, x_2, x_3 \geq 0
\end{align*}

\noindent
Since this LP is already in the standard form, we can construct the dual straightforward using methods discussed in class:
\begin{align*}
  &\mathrm{min} \quad W = 60y_1 + 10y_2 + 20y_3 \\[2pt]
  &\mathrm{s.t.} \quad 
  \begin{array}[t]{r c r c r c r}
      3y_1 & + & y_2 & - & y_3 & \geq & 2 \\
      y_1 & - & y_2 & + & y_3 & \geq & -1 \\
      y_1 & + & 2y_2 & - & y_3 & \geq & 3
  \end{array}\\
  &\mathrm{and} \quad y_1, y_2, y_3 \geq 0
\end{align*}

\subsection{B. Construct a dual using the SOB method.}
\begin{align*}
  &\mathrm{min} \quad Z = 4x_1 + 5x_2 + 3x_3 \\[2pt]
  &\mathrm{s.t.} \quad 
  \begin{array}[t]{r c r c r c r}
      x_1 & + & x_2 & + & 2x_3 & \geq & 20 \\
      -15x_1 & + & 6x_2 & - & 5x_3 & \leq & 50 \\
      x_1 & + & 3x_2 & + & 5x_3 & = & 30
  \end{array}\\
  &\mathrm{and} \quad x_1 \geq 0, x_2 \text{ unconstrained}, x_3 \leq 0
\end{align*}

\begin{itemize}
  \item The first constraint ($\geq$) is a \textit{sensible} constraint. This translates to $y_1 \geq 0$.
  \item The second constraint ($\leq$) is a \textit{bizarre} constraint. This translates to $y_2^\prime \leq 0$.
  \item The third constraint ($=$) is an \textit{odd} constraint. This translates to $y_3^\prime \text{ unconstrained}$.
\end{itemize}
\vspace{1mm}
\begin{itemize}
  \item The first variable ($x_1 \geq 0$) is \textit{sensible}. This translates to $y_1 - 15y_2^\prime + y_3^\prime \leq 4$.
  \item The second variable ($x_2$ unconstrained) is \textit{odd}. This translates to $y_1 + 6y_2^\prime + 3y_3^\prime = 5$.
  \item The third variable ($x_3 \leq 0$) is \textit{bizarre}. This translates to $2y_1 - 5y_2^\prime + 5y_3^\prime \geq 3$.
\end{itemize}

\noindent
Combining all of this gives the following dual:

\newpage
\begin{align*}
  &\mathrm{max} \quad W = 20y_1 + 50y_2^\prime + 30y_3^\prime \\[2pt]
  &\mathrm{s.t.} \quad 
  \begin{array}[t]{r c r c r c r}
      y_1 & - & 15y_2^\prime & + & y_3^\prime & \leq & 4 \\
      y_1 & + & 6y_2^\prime & + & 3y_3^\prime & = & 5 \\
      2y_1 & - & 5y_2^\prime & + & 5y_3^\prime & \geq & 3
  \end{array}\\
  &\mathrm{and} \quad y_1 \geq 0, y_2^\prime \leq 0, y_3^\prime \text{ unconstrained}
\end{align*}

%
% Exercise 3
%
\section{Alternative Definition of BFS}
\subsection{A. If $F \neq \varnothing$, there exists $C \subseteq F$ s.t. $C$ forms a cycle.}
Assume that if $F \neq \varnothing$, cycle $C \subseteq F$ does \underline{not} exist.

\noindent
\\
This means that there exists a fractional decision variable that connects a student with a task, but either the student or the task has no other fractional decision variable (because there is no cycle, the chain of fractional decision variables must end somewhere).

\noindent
\\
As a result of this, there exists at least one student or task whose sum of decision variables is a fractional value.

\noindent
\\
Since 1 is not a fractional value, this contradicts the constraints of the LP.

\noindent
\\
This means our original assumption that if $F \neq \varnothing$, cycle $C \subseteq F$ does not exist is false.
This proves that if $F \neq \varnothing$, there exists $C \subseteq F$ such that $C$ forms a cycle. $\Box$

\subsection{B. Adding $\epsilon$ to a decision variable $x \in C$.}
To regain feasibility, we can apply the following algorithm:

\noindent
\\
The idea is to cycle through all the decision variables $x \in C$, starting with one adjacent to the decision variable we increased by $\epsilon$, and ending once all decision variables $x \in C$ have an $\epsilon$ term (i.e. we have completed the cycle).
For each variable we do the following:

\begin{enumerate}
  \item If we increased the previous variable by $\epsilon$, decrease the current variable by $\epsilon$.
  \item If we decreased the previous variable by $\epsilon$, increase the current variable by $\epsilon$.
\end{enumerate}

\noindent
After this all the effects of adding $\epsilon$ to one of the variables will have been canceled out, and the solution will be feasible again.

\noindent
\\
For the non-negativity constraints to be satisfied, $\epsilon$ has to be greater than or equal to 0, and smaller than or equal to the smallest $x \in C$.

\subsection{C. Subtracting $\epsilon$ to a decision variable $x \in C$.}
For this we can use the same algorithm as described in the previous question (except this time step 2 will be called before step 1).

\newpage
\subsection{D. Prove that a BFS can never have fractional values.}
Let $x^*$ be a feasible solution with fractional values ($F \neq \varnothing$).
Our proof in question A tells us that $x^*$'s bipartite graph representation has a cycle $C \subseteq F$. 

\noindent
\\
Our solutions for questions B and C tell us that if we have a feasible solution with a cycle, we can add a value $\epsilon$ to one of the variables $x \in C$, and we can turn it into a feasible solution again (assuming $0 \leq \epsilon \leq \text{ min. }x \in C$).

\noindent
\\
Let $a = \frac{1}{2} (\text{min. } x \in C)$.

\noindent
\\
Let $x^\prime$ be $x^*$ with $a$ added to one of the variables $x \in C$.\\
Let $x^{\prime\prime}$ be $x^*$ with $a$ subtracted from one of the variables $x \in C$.\\
We know these solutions are feasible, because $0 \leq a \leq \text{min. } x \in C$.

\noindent
\\
Because these actions are each other's opposites, we can say that $x^*$ lies on the line segment constructed from $x^\prime$ to $x^{\prime\prime}$:
\[ x^* = 0.5 x^\prime + 0.5 x^{\prime\prime} \]

\noindent
This contradicts the line-segment definition of a BFS, which means $x^*$ cannot be a BFS.\\
This proves that a feasible solution with fractional values cannot be a BFS, which proves that all BFS's must be integral. $\Box$

\end{document}
